% Copyright (c) 2026 Mark Hunter. SPDX-License-Identifier: GPL-3.0-or-later
\chapter{Authority, identity and use}\label{guide}

This publication answers: ``What capability architecture does the governed RAD currently contain?''
Version 0.1 represents the current governed Markdown architecture without extending,
completing or reconciling it. The Markdown RAD remains architectural authority;
this publication is its definitive LaTeX representation at the recorded baseline.
The Architectural Axioms remain the highest design authority.

AX-01 and AX-04 preserve the business and health-information meaning of the model.
AX-14 preserves distinctions between capability, behaviour, exposure, progression,
relevance and status. AX-17 requires established, candidate and unresolved assertions
to remain distinguishable. AX-18 prevents integration/coordination from conferring
domain authority. This publication applies those constraints to representation;
it makes no architectural change. The selected axiom text, metamodel and K9
decision are reproduced in the semantic-detail part.

\section{Reading the model}

Business Capabilities use only the explicitly established CT1 → CT2 → CT3
hierarchy. Regions are contextual groupings, not ancestors. Business Enabling
Capabilities require CT1 → CT2 → CT3 → FT for a fully reconciled Feature, but the
current RAD does not establish those BEC tier allocations or complete ancestry.
This publication preserves established owners and the explicit partial
Client Administration → Person Identity responsibility chain without manufacturing
missing tiers. It reports zero \emph{established} BEC CT1, CT2 or CT3 allocations;
this does not mean that the 76 established BEC contexts do not exist.

Functions, Services and Processes retain their evidenced owners. They form a graph,
not a mandatory Feature → Function → Service/Process tree. A Function need not
have a Feature, Service or Process association. Enterprise Capabilities retain the
fourteen governed entries and their actual descriptive responsibilities; no additional
Enterprise decomposition is introduced.

Each element has a stable label, metadata and a link to exact RAD source detail.
The semantic-detail part preserves the source definitions, boundaries, state-change
statements, illustrative activities, examples, exclusions, notes, classifications,
relationship tables and unresolved matters. Its source headings and retained diagram
notation are not new containment assertions. The complete \texttt{manifest.json}
records exact source lines and fingerprints; \texttt{source-coverage.json} records the
transcribed blocks. Source references such as ``business:152'' mean source-key and
Markdown line, not a new architectural identifier.

\section{Definitive publication identity}

Existing Feature identifiers (for example \texttt{FEAT-EM-01}) and Enterprise
identifiers (\texttt{EC-01} through \texttt{EC-14}) are preserved exactly.
Catalogue ordinals, BC-18-style source shorthand and CT syntax examples are not
allocated structural Canonical IDs. No structural Canonical ID is allocated here.

For elements without an allocated identifier, the frozen Reference Id uses
\texttt{HCM-}\emph{type}\texttt{-}\emph{semantic-name-token}. Type tokens are
\texttt{BC}, \texttt{BEC}, \texttt{FN}, \texttt{SV} and \texttt{PR}.
The initial token is the lowercase name with runs of non-alphanumeric characters
replaced by a hyphen. This is a deterministic initial allocation rule only:
\textbf{after v0.1, the manifest value must be reused verbatim, including after
rename or relocation}. It is not regenerated from a future name, parent, tier,
ordinal, chapter or build. A collision requires an explicitly recorded allocation,
never a position-based suffix; this baseline has no collisions.

The Reference Id is a definitive publication identity, distinct from the metamodel's
contextual structural Canonical ID. The metamodel requires a structural ID to change
when responsibility context changes, retaining the former value as an alias.
Keeping these two fields separate satisfies both that rule and publication identity
stability. Unallocated or ancestry-dependent structural IDs remain unresolved.

Stable LaTeX label namespaces are \texttt{cap:bc-}, \texttt{cap:bec-},
\texttt{cap:ec-}, \texttt{feature:}, \texttt{function:}, \texttt{service:}
and \texttt{process:}. The exact values are frozen in the manifest and are also
preserved across relocation or rename. Source-detail labels are locators, not
architectural element identities. Internal reference-candidate record keys are
retained for navigation continuity only; they are neither Canonical IDs nor aliases.

\section{Architectural status and uncertainty}

\textbf{Canonical} means that the element's identity/responsibility is established
by the cited governed RAD. It does not establish every possible property or
relationship of that element. \textbf{Candidate} preserves a named provisional
concern without approving its placement or decomposition. \textbf{Unresolved}
qualifies a particular unestablished dimension or relationship. It is not a fifth
Harmonia relevance classification. Source-specific provisional proposals remain
provisional in their source prose; occurrence in a canonical document does not
promote them to Canonical.

Health Service Governance is a Canonical CT1 with unresolved relevance and enabling
derivation. The three assurance BECs are Canonical responsibilities with unresolved
tiers, ancestry, view placement, relevance and Feature decomposition. All 137
Features are established FT identities with unresolved complete capability ancestry.
The 69 enabling catalogue descriptions with unestablished standalone element
identity are not counted as new Capabilities.

\section{Reconciled publication counts}

\begin{longtable}{|p{0.53\linewidth}|p{0.37\linewidth}|}
\hline\rowcolor{TableHead}\textbf{Architectural category} & \textbf{RAD / publication count}\\\hline
Business CT1 / CT2 / CT3 & 18 / 39 / 133\\\hline
Canonical Business Capabilities & 190\\\hline
Named candidate BC concerns with unresolved placement & 3 (separate from 190)\\\hline
Established BEC CT1 / CT2 / CT3 assignments & 0 / 0 / 0; tiers unresolved\\\hline
Canonical BEC responsibility contexts & 76\\\hline
Features / Functions / Services / Processes & 137 / 147 / 133 / 19\\\hline
Enterprise Capabilities & 14\\\hline
Named enabling decomposition descriptions & 73; 4 independently established, 69 unresolved identities\\\hline
Definitive architectural-element publication records & 719; 716 Canonical, 3 Candidate\\\hline
\end{longtable}

Counts are of independently established identities or explicitly retained candidates,
not source occurrences, descriptive bullets or possible future peers. The reference
candidate capture was used to locate current Markdown evidence; its candidate
standing supplies no architectural approval. Counts and identifiers are checked
against the actual Markdown inventories and explicit parent assertions.

\section{Preserved issues requiring later architectural review}

The full source qualifications remain controlling. In particular:
\begin{itemize}
\item BEC tier/ancestry reconciliation, 137 Feature frames and 69 description
identities remain unresolved under K9; the five views do not resolve them.
\item Health Information Exchange identifies distinct BC and BEC responsibilities.
Their naming collision remains unresolved; they are not merged or renamed.
\item Medication compounding/preparation (OI-BC-001), operational-readiness
placement (OI-BC-002), maintenance membership and sharing-programme placement
remain unadjudicated. Clinical Intervention's proposed three-child taxonomy
remains provisional; Clinical Determination and Individual Care Planning have
no adjudicated CT3 set.
\item Operational Credential Management remains distinct from Clinical
Qualification Management and its historical name alias. Name similarity does
not establish Feature association.
\item The existing care-team collaboration Function description mentions managing
clinical handovers while the metamodel excludes that authority. Both source
assertions are retained for review; this publication does not resolve the conflict.
\item Process illustrations retain their stated scope and unresolved transitions,
including Encounter ON\_LEAVE and the disconnected stalled/failed Synthetic Task
outcome. Assurance lifecycle detail and internal Service/Function mappings remain
unestablished; no additional state or edge is inferred.
\item The assurance Strategy source's opening reconciliation qualifies older BC-18
tier/relevance wording. Its Strategy stop boundary does not erase the later
approved Domain03 assurance catalogue. B2.1's retained L1/L2 terminology does not
allocate or override current CT classifications.
\end{itemize}

No legacy documentation or architectural material from the earlier candidate LaTeX
publications was used. Version 0.2 derivation and reconciliation are outside this baseline.
